\section{Post-Estimation Infrastructure}
\label{sec:post_estimation}

EconIRL provides four methods for computing standard errors on estimated parameters. Asymptotic standard errors use the inverse of the negative Hessian, $\mathrm{Var}(\hat\theta) = (-H)^{-1}$, which is valid under correct specification of the likelihood. Robust sandwich standard errors account for potential misspecification by computing $H^{-1} B H^{-1}$, where $B$ is the outer product of per-observation score vectors. Nonparametric bootstrap standard errors resample trajectories with replacement and re-estimate the model on each bootstrap sample, yielding confidence intervals that do not rely on asymptotic normality. Clustered standard errors aggregate the score vectors within individuals before forming the outer product, accounting for within-individual serial correlation in panel data. The Hessian is computed analytically for NFXP via implicit differentiation through the Bellman fixed point following \citet{iskhakov2016}, and numerically via central finite differences for all other estimators.

The library runs four identification diagnostics before estimation begins. Feature matrix rank is computed on the flattened feature matrix of dimension $(|\mathcal{S}| \times |\mathcal{A}|) \times K$, and if the rank falls below the number of parameters $K$, the model is under-identified and estimation is halted. The condition number of the feature matrix flags near-collinearity that inflates standard errors even when the matrix is technically full rank, with values above $10^6$ triggering a warning. State coverage reports how many of the $|\mathcal{S}|$ states have at least one observation in the data, because states with zero observations produce degenerate CCP estimates that cannot inform the likelihood. Single-action state counts identify states where only one action is ever observed, which produce boundary CCPs of zero or one and yield uninformative Bellman updates. After estimation completes, eigenvalue analysis of the negative Hessian verifies that the matrix is positive definite at the optimum, and the library reports any parameter correlations exceeding 0.95 as a warning of weak identification \citep{magnacThesmar2002}.

All estimators return an \texttt{EstimationSummary} object that contains parameter estimates, standard errors, $t$-statistics, $p$-values, 95 percent confidence intervals, and Wald test statistics for joint hypotheses. The \texttt{.summary()} method prints a formatted table in the style of the \texttt{statsmodels} library, with one row per parameter and columns for the point estimate, standard error, $t$-statistic, and confidence bounds. The \texttt{.to\_latex()} method produces a publication-ready table using the \texttt{booktabs} package, suitable for direct inclusion in manuscripts. This uniform output format means that switching from NFXP to MCE-IRL or from CCP to AIRL requires changing only the estimator class, with no changes to downstream inference or reporting code.

Given estimated parameters $\hat\theta$, the library supports counterfactual simulation by re-solving the Bellman equation under alternative parameters $\theta'$ and computing the welfare change $\Delta W(s) = V(s;\theta') - V(s;\hat\theta)$ for each state. This allows researchers to quantify the value of policy interventions such as a tax on replacement costs or a subsidy for maintenance. Elasticity analysis sweeps a single parameter over a user-specified range while holding all other parameters fixed, re-solves for the optimal policy at each point, and reports the resulting policy response surface. The counterfactual module also supports changes to the transition matrix, enabling analysis of how improvements in component reliability or changes in usage patterns affect optimal replacement timing and expected welfare.
